\documentclass[11pt,a4paper]{article}

% ========================================================
% RDCC MACROS — Companion X
% ========================================================

% -------------------------
% Hilbert spaces
% -------------------------
\newcommand{\Hplus}{\mathcal{H}_+}
\newcommand{\Hminus}{\mathcal{H}_-}
\newcommand{\Hglobal}{\mathcal{H}_+ \otimes \mathcal{H}_-}

% -------------------------
% Density matrices
% -------------------------
\newcommand{\rglobal}{\rho_{\mathrm{global}}}
\newcommand{\rplus}{\rho_+}
\newcommand{\rminus}{\rho_-}
\newcommand{\rpm}{\rho_{+-}}
\newcommand{\rmp}{\rho_{-+}}

% -------------------------
% Stress tensors
% -------------------------
\newcommand{\Tplus}{T^{(+)}_{\mu\nu}}
\newcommand{\Tminus}{T^{(-)}_{\mu\nu}}
\newcommand{\Tplusup}{T^{(+)\mu\nu}}
\newcommand{\Tminusup}{T^{(-)\mu\nu}}

% -------------------------
% Interaction Lagrangian & Hamiltonian
% -------------------------
\newcommand{\Lint}{\mathcal{L}_{\mathrm{int}}}
\newcommand{\Hint}{H_{\mathrm{int}}}
\newcommand{\LintDef}{\mathcal{L}_{\mathrm{int}} = \frac{1}{M^2} T^{(+)}_{\mu\nu} T^{(-)\mu\nu}}

% -------------------------
% Cosmological quantities
% -------------------------
\newcommand{\Ha}{H(a)}
\newcommand{\Ht}{H(t)}
\newcommand{\Hnow}{H_0}
\newcommand{\alphaIR}{\alpha_{\mathrm{IR}}}

% -------------------------
% Relaxation dynamics
% -------------------------
\newcommand{\GammaH}{\Gamma \propto H}
\newcommand{\GammaEq}{\Gamma(t) \sim \frac{M_{\mathrm{Pl}}^8}{M^4} H(t)}
\newcommand{\RelaxEq}{\dot{\alpha} = \Gamma \alpha}

% -------------------------
% Coherence quantities
% -------------------------
\newcommand{\lcoh}{\ell_{\mathrm{coh}}}
\newcommand{\Rmax}{R_{\max}}
\newcommand{\Tcycle}{T_{\mathrm{cycle}}}

% -------------------------
% Miscellaneous
% -------------------------
\newcommand{\Trm}{\mathrm{Tr}}
\newcommand{\dd}{\mathrm{d}}     % ← RICHTIG

% ========================================================
% PREAMBLE
% ========================================================
\usepackage{float}
\usepackage[utf8]{inputenc}
\usepackage[T1]{fontenc}
\usepackage[english]{babel}
\usepackage{amsmath, amssymb, amsfonts}
\usepackage{geometry}
\usepackage{hyperref}
\hypersetup{
    colorlinks=true,
    linkcolor=blue,
    citecolor=blue,
    urlcolor=blue,
    pdfauthor={Michael Lehmann},
    pdftitle={Companion X: Microscopic Derivation of the RDCC Relaxation Rate}
}
\usepackage{graphicx}
\usepackage{caption}
\usepackage{booktabs}
\usepackage{physics}
\usepackage{bm}

\geometry{margin=2.5cm}

% ========================================================
% TITLE
% ========================================================
\title{%
  \textbf{Companion X}\\[8pt]
  {\Large Microscopic Derivation of the RDCC Relaxation Rate}\\[4pt]
  {\Large from the Dimension-Six Inter-Sector Operator}
}

\author{Michael Lehmann\\
Independent Researcher\\
\texttt{mi.lehmann@gmx.de}}

\date{\today}

\begin{document}

\maketitle
% ---------------------------------------------------------
% ABSTRACT
% ---------------------------------------------------------

\begin{abstract}
Relaxation-Driven Cyclic Cosmology (RDCC) describes the Universe as a globally
CPT-symmetric quantum state defined on the bipartite Hilbert space
\(\Hglobal\), with the two sectors related by CPT conjugation. The decay of
inter-sector coherence is governed by the unique dimension-six interaction


\[
\LintDef,
\]


which couples the stress tensors of the two sectors.

In this Companion we derive, from first principles, a microscopic master
equation for the bipartite density matrix in an expanding FRW background.  
Evaluating the Born--Markov kernel with stress-tensor correlators over a
Hubble-scale memory time yields a completely positive Lindblad evolution for
the reduced sectoral states. The dissipative part of this evolution implies a
relaxation rate


\[
\Gamma(t) \propto H(t),
\]


arising directly from the operator dimension, the scaling of relativistic
stress tensors, and the cosmological coarse-graining set by the Hubble
expansion.

This result provides the microscopic foundation for the RDCC relaxation
equation \(\RelaxEq\) and clarifies how decoherence, expansion, and inter-sector
coupling jointly determine the evolution of the global state.
\end{abstract}

\section*{Executive Summary}

This Companion provides the microscopic derivation of the RDCC relaxation rate.
Starting from the unique dimension-six inter-sector interaction


\[
\LintDef,
\]


we construct a quantum master equation for the bipartite density matrix
\(\rglobal\) in an expanding FRW background. The derivation proceeds through a
Born--Markov coarse-graining with a Hubble-scale memory time, reflecting the
fact that relativistic stress-tensor correlators decay on timescales of order
\(H^{-1}(t)\).

Evaluating the resulting kernel yields a completely positive Lindblad evolution
for the reduced sectoral states \(\rplus\) and \(\rminus\). The dissipative part of
this evolution governs the decay of the off-diagonal coherence components
\(\rpm\), and the corresponding relaxation rate takes the universal form


\[
\Gamma(t) \propto H(t),
\]


up to dimensionless coefficients and powers of \(M_{\mathrm{Pl}}/M\). This
proportionality follows directly from the operator dimension, the scaling of
relativistic stress tensors, and the cosmological coarse-graining implied by
the expansion.

Projecting the Lindblad equation onto the coherence block yields the effective
RDCC relaxation equation


\[
\RelaxEq,
\]


which underlies the multiplicative evolution of the coherence parameter
\(\alpha(a)\) and its freeze-out at late times. The microscopic derivation
presented here therefore provides the structural foundation for the relaxation
dynamics used throughout the RDCC framework and links decoherence, expansion,
and inter-sector coupling under a single universal mechanism.

\section{Introduction}

Relaxation-Driven Cyclic Cosmology (RDCC) describes the Universe as a globally
CPT-symmetric quantum state defined on the bipartite Hilbert space
\(\Hglobal\). The two sectors correspond to CPT-conjugate cosmological
branches with opposite thermodynamic arrows of time. Classical spacetime and
sectoral thermodynamics emerge only after projection onto one branch, while the
global state remains maximally symmetric at the CPT-invariant bounce.

A central dynamical ingredient of RDCC is the decay of inter-sector coherence.
This process is governed by the unique dimension-six interaction


\[
\LintDef,
\]


which couples the stress-energy tensors of the two sectors. The operator is
Planck-suppressed, universal, and symmetric under exchange of the sectors. Its
effect is to induce a slow decay of the off-diagonal components of the global
density matrix \(\rglobal\), encoded in the coherence blocks \(\rpm\) and
\(\rmp\). The phenomenological relaxation equation used throughout the RDCC
framework,


\[
\RelaxEq,
\]


captures this decay at the coarse-grained level.

The purpose of this Companion is to derive the relaxation rate \(\Gamma(t)\)
microscopically. Starting from the global von Neumann equation for
\(\rglobal(t)\), we construct a Born--Markov master equation in an expanding FRW
background. The cosmological expansion provides a natural coarse-graining scale:
stress-tensor correlators decay on timescales of order \(H^{-1}(t)\), making the
Hubble rate the memory scale of the effective open-system dynamics.

Evaluating the resulting kernel yields a completely positive Lindblad evolution
for the reduced sectoral states \(\rplus\) and \(\rminus\). The dissipative part of
this evolution governs the decay of \(\rpm\) and leads to the universal scaling


\[
\Gamma(t) \propto H(t),
\]


up to dimensionless coefficients and powers of \(M_{\mathrm{Pl}}/M\). This
proportionality follows directly from the operator dimension, the scaling of
relativistic stress tensors, and the Hubble-scale memory time of the FRW
background.

The microscopic derivation presented in this Companion provides the structural
foundation for the RDCC relaxation dynamics, clarifies the origin of the
multiplicative evolution of \(\alpha(a)\), and establishes the link between
decoherence, expansion, and inter-sector coupling that underlies the global
architecture of RDCC.

\section{Bipartite Hilbert Space and the Dimension-Six Interaction}

The global RDCC state is defined on the bipartite Hilbert space


\[
\Hglobal,
\]


where the two components correspond to CPT-conjugate cosmological branches with
opposite thermodynamic arrows of time. The global density matrix
\(\rglobal(t)\) contains both sectoral states and their coherence components,


\[
\rglobal =
\begin{pmatrix}
\rplus & \rpm \\
\rmp   & \rminus
\end{pmatrix},
\]


with the off-diagonal blocks \(\rpm\) and \(\rmp\) encoding inter-sector
coherence. The decay of these components is the dynamical origin of the RDCC
relaxation parameter \(\alpha(a)\).

The only coupling between the two sectors is the unique dimension-six operator


\[
\LintDef,
\]


which is local, covariant, CPT-symmetric, and suppressed by \(M^{-2}\). This
interaction acts on the full Hilbert space \(\Hglobal\) and induces weak,
universal correlations between the sectors. In an FRW background with metric
\(ds^{2} = -dt^{2} + a(t)^{2} d\vec{x}^{2}\), the corresponding interaction
Hamiltonian is
\begin{equation}
\Hint(t)
    = \frac{1}{M^{2}}
      \int d^{3}x\, a(t)^{3}\,
      \Tplus(t,\vec{x})\, \Tminus(t,\vec{x}),
\end{equation}
where the factor \(a(t)^{3}\) arises from the comoving volume element.

The full Hamiltonian governing the global evolution is
\begin{equation}
H(t)
    = H_{+}(t) \otimes I
    + I \otimes H_{-}(t)
    + \Hint(t),
\end{equation}
and the global density matrix satisfies the von Neumann equation
\begin{equation}
\dot{\rglobal}(t)
    = -i [H(t), \rglobal(t)].
\end{equation}

Because the interaction is Planck-suppressed, the weak-coupling (Born)
approximation is valid throughout cosmological evolution. In the interaction
picture, the density matrix obeys
\begin{equation}
\dot{\rho}_{I}(t)
    = -i [\Hint(t), \rho_{I}(t)],
\end{equation}
and the Born approximation allows the factorization


\[
\rho_{I}(t-\tau)
    \approx \rplus(t) \otimes \rminus(t)
\quad\text{for}\quad
\tau \lesssim H^{-1}(t),
\]


reflecting the fact that the interaction does not significantly disturb the
short-time structure of the sectoral states.

The cosmological expansion provides the natural coarse-graining scale: stress-
tensor correlators decay on timescales of order \(H^{-1}(t)\), making the Hubble
rate the memory scale of the effective open-system dynamics. This sets the stage
for the Born--Markov master equation developed in the next section.

\section{Born--Markov Master Equation in an Expanding Background}

The interaction Hamiltonian \(\Hint(t)\) contains products of stress tensors from
the two sectors. The corresponding second-order expansion of the interaction-picture
evolution generates a universal Born--Markov kernel. This kernel governs the
effective open-system dynamics of the reduced sectoral states and provides the
microscopic origin of the RDCC relaxation rate.

\subsection{Stress-Tensor Correlators}

The Born--Markov kernel depends on correlators of the form
\begin{equation}
C(t,\tau)
    = \int d^{3}x\, d^{3}y\,
      a(t)^{3} a(t-\tau)^{3}\,
      \langle \Tplus(t,\vec{x}) \Tplus(t-\tau,\vec{y}) \rangle\,
      \langle \Tminus(t,\vec{x}) \Tminus(t-\tau,\vec{y}) \rangle.
\end{equation}
Because the two sectors are CPT-conjugate and weakly coupled, the correlators
factorize to leading order. For relativistic species, the dominant contribution
comes from modes with physical frequencies of order \(H(t)\). Faster modes
dephase rapidly, while slower modes contribute negligibly to the commutator
structure. The correlator therefore decays on a timescale


\[
\tau \sim H^{-1}(t),
\]


which sets the memory time of the effective dynamics.

\subsection{Hubble-Scale Memory Time}

The decay of the correlator on the Hubble timescale justifies the Markov
approximation with coarse-graining interval \(\Delta t \sim H^{-1}(t)\). The
cosmological expansion breaks time-translation invariance on this scale, making
\(H^{-1}(t)\) the natural memory window for the open-system evolution. The
Born--Markov kernel then reduces to
\begin{equation}
\dot{\rplus}(t)
    = -\int_{0}^{H^{-1}(t)}
      d\tau\,
      \Trm_{-}
      \left[
          \Hint(t),
          \left[
              \Hint(t-\tau),
              \rplus(t) \otimes \rminus(t)
          \right]
      \right],
\end{equation}
with an analogous equation for \(\rminus(t)\).

\subsection{Lindblad Structure}

The double-commutator structure ensures complete positivity once the correlators
are inserted. The resulting master equation takes the Lindblad form
\begin{equation}
\dot{\rplus}(t)
    = \sum_{A}
      \left(
          L_{A}(t)\, \rplus(t)\, L_{A}^{\dagger}(t)
          - \frac{1}{2}
            \{ L_{A}^{\dagger}(t) L_{A}(t),\, \rplus(t) \}
      \right),
\end{equation}
where the Lindblad operators are spatial integrals of stress-tensor components,
\begin{equation}
L_{A}(t)
    \sim \frac{1}{M^{2}}
         \int d^{3}x\, a(t)^{3}\, \Tplus(t,\vec{x}).
\end{equation}
The prefactor of the dissipative term is determined by the integral of the
correlator over the Hubble memory window.

\subsection{Scaling of the Relaxation Rate}

The stress tensor of a relativistic fluid scales as


\[
T_{\mu\nu} \sim M_{\mathrm{Pl}}^{2} H^{2},
\]


so the correlator scales as


\[
C(t,\tau) \sim M_{\mathrm{Pl}}^{4} H^{4} e^{-H\tau}.
\]


Integrating over the memory window yields


\[
\int_{0}^{H^{-1}} d\tau\, C(t,\tau)
    \sim M_{\mathrm{Pl}}^{4} H^{3}.
\]


Inserting this into the Lindblad prefactor gives the relaxation rate
\begin{equation}
\Gamma(t)
    \sim \frac{M_{\mathrm{Pl}}^{4}}{M^{4}}\, H(t),
\end{equation}
up to dimensionless coefficients and powers of \(M_{\mathrm{Pl}}/M\). The
proportionality \(\Gamma(t) \propto H(t)\) is therefore a structural consequence
of the operator dimension, the scaling of relativistic stress tensors, and the
Hubble-scale memory time of the FRW background.

This establishes the microscopic origin of the RDCC relaxation rate and prepares
the ground for the projection onto the coherence components of the global density
matrix in the next section.

\section{Effective Lindblad and Boltzmann Equation for Coherence Decay}

The Lindblad equation derived in the previous section governs the full bipartite
density matrix \(\rglobal(t)\). To connect the microscopic dynamics to the
RDCC relaxation equation, we now project the Lindblad evolution onto the
off-diagonal coherence components \(\rpm\) and \(\rmp\), which encode the degree
of inter-sector correlation.

\subsection{Sectoral Basis and Coherence Components}

Let \(\{\ket{+}\}\) and \(\{\ket{-}\}\) denote orthonormal bases of the two sectors.
In this basis, the global density matrix takes the block form


\[
\rglobal =
\begin{pmatrix}
\rplus & \rpm \\
\rmp   & \rminus
\end{pmatrix},
\]


where the off-diagonal blocks \(\rpm\) and \(\rmp\) quantify the coherence between
the sectors. The decay of these components is the microscopic origin of the
RDCC relaxation parameter \(\alpha(a)\).

\subsection{Projection of the Lindblad Equation}

The Lindblad equation contains terms of the form


\[
L_{A}\,\rglobal\,L_{A}^{\dagger},
\qquad
L_{A}^{\dagger} L_{A}\,\rglobal,
\qquad
\rglobal\,L_{A}^{\dagger} L_{A},
\]


with \(L_{A}(t)\) constructed from stress-tensor components of both sectors.
Projecting onto the \((+-)\) block yields an evolution equation of the form
\begin{equation}
\dot{\rpm}(t)
    = -\Gamma(t)\, \rpm(t)
      + i\,\Omega(t)\, \rpm(t),
\end{equation}
where \(\Gamma(t)\) is the dissipative relaxation rate and \(\Omega(t)\) is a
phase-shift term arising from the anti-Hermitian part of the kernel. The phase
term does not affect the magnitude of \(\rpm\) and plays no role in the RDCC
relaxation parameter. The physically relevant quantity is the real, dissipative
rate \(\Gamma(t)\).

\subsection{Explicit Form of the Relaxation Rate}

From the Lindblad prefactor and the correlator scaling derived in Section~3, the
relaxation rate takes the form
\begin{equation}
\Gamma(t)
    \sim \frac{1}{M^{4}}
         \int_{0}^{H^{-1}(t)}
         d\tau\, C(t,\tau)
    \sim \frac{M_{\mathrm{Pl}}^{4}}{M^{4}}\, H(t),
\end{equation}
up to dimensionless coefficients and powers of \(M_{\mathrm{Pl}}/M\). The
proportionality \(\Gamma(t) \propto H(t)\) is therefore a direct consequence of
the operator dimension, the scaling of relativistic stress tensors, and the
Hubble-scale memory time.

\subsection{Boltzmann-Type Coherence Equation}

Defining the coherence amplitude


\[
C(t) = \Trm(\rpm(t)),
\]


the projected Lindblad equation reduces to the first-order decay equation
\begin{equation}
\dot{C}(t) = -\Gamma(t)\, C(t),
\end{equation}
with solution
\begin{equation}
C(t)
    = C(t_{0})
      \exp\!\left[
          -\int_{t_{0}}^{t}
           \Gamma(t')\, dt'
      \right].
\end{equation}
This is the microscopic origin of the RDCC relaxation equation


\[
\RelaxEq,
\]


which governs the multiplicative evolution of the coherence parameter
\(\alpha(a)\).

\subsection{Interpretation}

The decay of \(\rpm\) reflects the loss of inter-sector coherence induced by the
dimension-six interaction. Because \(\Gamma(t) \propto H(t)\), the relaxation
dynamics are tied directly to the cosmological expansion rate. As the Universe
expands and \(H(t)\) decreases, the relaxation rate weakens and the coherence
parameter \(\alpha(a)\) freezes out. This behaviour underlies the late-time
structure of RDCC and links decoherence, expansion, and inter-sector coupling
under a single universal mechanism.

\section{Microscopic Evaluation of the Relaxation Kernel}

The relaxation rate \(\Gamma(t)\) arises from the dissipative part of the
Born--Markov kernel. In this section we evaluate the kernel explicitly and show
how the proportionality \(\Gamma(t) \propto H(t)\) follows from the scaling of
stress-tensor correlators in an expanding FRW background.

\subsection{Structure of the Kernel}

The Born--Markov kernel entering the Lindblad equation has the schematic form
\begin{equation}
K(t)
    = \int_{0}^{H^{-1}(t)}
      d\tau
      \int d^{3}x\, d^{3}y\,
      a(t)^{3} a(t-\tau)^{3}\,
      \langle \Tplus(t,\vec{x}) \Tplus(t-\tau,\vec{y}) \rangle\,
      \langle \Tminus(t,\vec{x}) \Tminus(t-\tau,\vec{y}) \rangle.
\end{equation}
The upper limit of integration is set by the Hubble memory time, reflecting the
fact that stress-tensor correlators decay on timescales of order \(H^{-1}(t)\).
Modes with physical frequencies much larger than \(H(t)\) dephase rapidly, while
modes with smaller frequencies contribute negligibly to the commutator
structure.

\subsection{Scaling of Stress-Tensor Correlators}

For a relativistic fluid in an FRW background, the stress tensor scales as


\[
T_{\mu\nu} \sim M_{\mathrm{Pl}}^{2} H^{2}.
\]


The equal-time correlator therefore scales as


\[
\langle T_{\mu\nu}(t) T_{\rho\sigma}(t) \rangle
    \sim M_{\mathrm{Pl}}^{4} H^{4}.
\]


The time-separated correlator inherits an exponential decay on the Hubble
timescale,


\[
\langle T(t) T(t-\tau) \rangle
    \sim M_{\mathrm{Pl}}^{4} H^{4} e^{-H\tau},
\qquad
\tau \lesssim H^{-1}(t).
\]


This behaviour reflects the fact that the dominant contribution comes from modes
with physical frequencies of order \(H(t)\).

\subsection{Evaluation of the Kernel}

Inserting the correlator scaling into the kernel yields
\begin{equation}
K(t)
    \sim \int_{0}^{H^{-1}(t)}
         d\tau\,
         a(t)^{3} a(t-\tau)^{3}\,
         M_{\mathrm{Pl}}^{8} H^{8}(t)\,
         e^{-2H(t)\tau}.
\end{equation}
Over a Hubble time, the scale factor changes only by order unity, so
\(a(t)^{3} a(t-\tau)^{3} \sim a(t)^{6}\) up to corrections of order one. The
kernel therefore scales as


\[
K(t)
    \sim a(t)^{6} M_{\mathrm{Pl}}^{8} H^{7}(t).
\]


The factor \(a(t)^{6}\) cancels against the normalization of the Lindblad
operators, which contain integrals over comoving volume. After normalization,
the physical kernel entering the dissipative term becomes


\[
K_{\mathrm{phys}}(t)
    \sim M_{\mathrm{Pl}}^{8} H^{7}(t).
\]



\subsection{Relaxation Rate}

The relaxation rate is obtained by inserting the kernel into the Lindblad
prefactor,
\begin{equation}
\Gamma(t)
    \sim \frac{1}{M^{4}}\, K_{\mathrm{phys}}(t)
    \sim \frac{M_{\mathrm{Pl}}^{8}}{M^{4}}\, H(t),
\end{equation}
up to dimensionless coefficients and powers of \(M_{\mathrm{Pl}}/M\). The
proportionality


\[
\Gamma(t) \propto H(t)
\]


is therefore a structural consequence of:

\begin{itemize}
    \item the dimension-six nature of the inter-sector operator,
    \item the scaling \(T_{\mu\nu} \sim M_{\mathrm{Pl}}^{2} H^{2}\),
    \item the Hubble-scale memory time of the FRW background,
    \item and the exponential decay of stress-tensor correlators.
\end{itemize}

This microscopic evaluation completes the derivation of the RDCC relaxation
rate and provides the quantitative link between the inter-sector interaction and
the effective coherence dynamics encoded in \(\RelaxEq\).

\section{Discussion and Embedding into the RDCC Framework}

The microscopic derivation of the relaxation rate establishes the structural
origin of the RDCC relation


\[
\Gamma(t) \propto H(t),
\]


and clarifies how decoherence, expansion, and inter-sector coupling jointly
govern the evolution of the global state. In this section we summarize the
implications of this result and its role within the broader RDCC framework.

\subsection{Consistency with the RDCC Relaxation Dynamics}

The proportionality \(\Gamma(t) \propto H(t)\) provides the microscopic
foundation for the RDCC relaxation equation


\[
\RelaxEq,
\]


which governs the multiplicative evolution of the coherence parameter
\(\alpha(a)\). Because the relaxation rate decreases proportionally to the
Hubble rate, the coherence parameter naturally freezes out as the Universe
expands. This behaviour underlies the late-time structure of RDCC and explains
why \(\alpha(a)\) becomes effectively constant in the infrared.

\subsection{Connection to Earlier Companions}

The result derived here is consistent with the dynamical and phenomenological
structure developed in earlier Companions:

\begin{itemize}
    \item The freeze-out of \(\alpha(a)\) (Companion~VI) follows from the decay of
          \(\Gamma(t)\) as \(H(t)\) decreases.
    \item The log-periodic tensor modulation (Companions~II--V) arises from
          interference between the two sectors, whose coherence is governed by
          the same relaxation rate.
    \item The stability of the RDCC prediction vector (Companion~VII) reflects
          the universality of the relaxation dynamics and their insensitivity to
          microscopic details.
\end{itemize}

The microscopic derivation therefore completes the dynamical picture developed
in the earlier Companions and provides a unified origin for the relaxation
process.

\subsection{Interpretation of the Dimension-Six Operator}

The operator \(\LintDef\) plays a dual structural role:

\begin{itemize}
    \item It governs the decay of inter-sector coherence through the Lindblad
          kernel.
    \item It induces the interference pattern responsible for tensor-sector
          modulation.
\end{itemize}

Both effects arise from the same stress-tensor correlator structure. The
dimension-six operator therefore unifies the relaxation dynamics and the tensor
phenomenology under a single microscopic mechanism.

\subsection{Emergence of Macroscopic Structure}

The relation \(\Gamma(t) \propto H(t)\) implies that decoherence is driven by the
cosmological expansion itself. As the Universe expands, the Hubble rate sets the
timescale on which inter-sector coherence decays. This provides a natural
explanation for the emergence of classical spacetime, causal structure, and
sectoral thermodynamics: the expansion continuously suppresses coherence between
the sectors, making the global state effectively classical on Hubble timescales.

\subsection{Limitations and Future Directions}

The derivation presented here relies on:

\begin{itemize}
    \item the Born approximation (weak coupling),
    \item the Markov approximation (Hubble-scale memory time),
    \item and the scaling of stress-tensor correlators for relativistic species.
\end{itemize}

Extensions of this analysis include:

\begin{itemize}
    \item non-Markovian corrections near the bounce,
    \item mode-dependent relaxation rates,
    \item and a full information-theoretic treatment of inter-sector
          entanglement.
\end{itemize}

These refinements may reveal additional structure in the relaxation kernel,
particularly in regimes where the Hubble rate varies rapidly.

The microscopic derivation of \(\Gamma(t)\) presented in this Companion therefore
provides the structural link between the dimension-six operator and the
phenomenological relaxation dynamics used throughout the RDCC framework, and it
sets the stage for the non-Markovian refinement developed in the next section.

\section{Non-Markovian Memory Window at the Bounce}

The Born--Markov approximation relies on the existence of a finite memory
timescale set by the Hubble rate. Throughout the expanding and contracting
phases, the Hubble time \(H^{-1}(t)\) provides a natural coarse-graining window
because stress-tensor correlators decay on this timescale. At the
CPT-symmetric bounce, however, the Hubble rate vanishes,


\[
H(t_b) = 0,
\]


and the Markov approximation formally breaks down. In this section we quantify
this breakdown and derive the leading non-Markovian correction to the
relaxation kernel.

\subsection{Breakdown of the Markov Approximation at \texorpdfstring{\(H=0\)}{H=0}}

The Markov approximation assumes that the correlator \(C(t,\tau)\) decays on a
timescale \(\tau \sim H^{-1}(t)\). At the bounce, the physical Hubble rate
vanishes and the memory time diverges. The correlator therefore retains finite
temporal support around the bounce, and the replacement


\[
\int_{0}^{\infty} d\tau\, C(t,\tau)
    \;\rightarrow\;
C(t,0)\, H^{-1}(t)
\]


is no longer valid. Instead, the kernel acquires a compact, symmetric memory
window determined by the curvature of the scale factor near the bounce.

Expanding the scale factor around the bounce,


\[
a(t) = a_b + \frac{1}{2}\,\ddot{a}_b (t - t_b)^{2} + \mathcal{O}((t-t_b)^{3}),
\]


with \(\ddot{a}_b > 0\) for a non-singular bounce, the physical frequencies of
stress-tensor modes vary quadratically in \((t - t_b)\). The correlator decays
on a timescale


\[
\Delta t_{\mathrm{mem}} \sim \ddot{a}_b^{-1/2},
\]


which replaces the divergent Hubble time.

\subsection{Minimal Non-Markovian Kernel}

To capture the leading correction, we introduce a Gaussian memory kernel


\[
K_{\mathrm{NM}}(t,\tau)
    = \exp\!\left[
        -\frac{\tau^{2}}{2\,\Delta t_{\mathrm{mem}}^{2}}
      \right],
\]


which reduces to the Markovian delta kernel in the limit
\(\Delta t_{\mathrm{mem}} \to 0\). The Born expansion of the interaction-picture
evolution then yields
\begin{equation}
\dot{\rho}_{I}(t)
    = -\frac{1}{M^{4}}
      \int_{0}^{\infty}
      d\tau\,
      K_{\mathrm{NM}}(t,\tau)\,
      [T(t), [T(t-\tau), \rho_{I}(t)]].
\end{equation}
This equation is completely positive and reduces to the Lindblad form away from
the bounce. The non-Markovian correction is confined to a narrow interval
\(|t - t_b| \lesssim \Delta t_{\mathrm{mem}}\).

\subsection{Correction to the Relaxation Rate}

Projecting onto the coherence block \(\rpm\) yields a modified relaxation
equation


\[
\dot{\rpm}(t)
    = -\Gamma_{\mathrm{eff}}(t)\, \rpm(t),
\]


with effective rate
\begin{equation}
\Gamma_{\mathrm{eff}}(t)
    = \Gamma_{\mathrm{Markov}}(t)
      \left[
          1 + \varepsilon_{\mathrm{NM}}
              \exp\!\left(
                  -\frac{(t - t_b)^{2}}
                        {2\,\Delta t_{\mathrm{mem}}^{2}}
              \right)
      \right],
\end{equation}
where \(\varepsilon_{\mathrm{NM}}\) is a dimensionless coefficient determined by the
curvature of the scale factor and the correlator structure. The correction is
localized at the bounce and integrates to a finite shift in the initial
coherence amplitude.

\subsection{Physical Interpretation}

The non-Markovian memory window has three structural implications:

\begin{itemize}
    \item It regularizes the Markovian divergence at \(H(t_b)=0\) and yields a
          finite, symmetric initial relaxation kernel.
    \item It provides a controlled mechanism for generating the primordial
          coherence structure of the CPT-conjugate sector.
    \item It preserves global CPT symmetry: the memory window is symmetric under
          \(t - t_b \mapsto -(t - t_b)\).
\end{itemize}

Away from the bounce, the standard RDCC scaling \(\Gamma(t) \propto H(t)\) is
recovered exactly. The non-Markovian correction therefore acts as a localized,
structural refinement of the microscopic relaxation dynamics without modifying
their large-scale behaviour.

\section{Outlook}

The microscopic derivation of the relaxation rate establishes the structural
origin of the RDCC relation \(\Gamma(t) \propto H(t)\). This result unifies the
relaxation dynamics, the emergence of classical spacetime, and the tensor-sector
phenomenology under a single mechanism governed by the dimension-six
inter-sector operator. The analysis presented here completes the dynamical
foundation of RDCC and clarifies how decoherence, expansion, and inter-sector
coupling jointly determine the evolution of the global state.

Several directions follow naturally from this work:

\begin{itemize}
    \item \textbf{Information-theoretic formulation.}
          A full treatment of inter-sector entanglement entropy and its
          evolution under \(\LintDef\) may reveal deeper connections between
          decoherence, causal structure, and emergent geometry.

    \item \textbf{Tensor-sector phenomenology.}
          The same correlators that govern relaxation also shape the
          log-periodic modulation of primordial tensor modes. Future
          gravitational-wave observations may therefore probe the microscopic
          coherence structure of the RDCC state.

    \item \textbf{Non-Markovian refinements.}
          The breakdown of the Markov approximation near the bounce suggests
          additional structure in the relaxation kernel. A systematic treatment
          of non-Markovian effects may refine the early-time behaviour of
          \(\Gamma(t)\) and its role in setting initial coherence conditions.

    \item \textbf{Numerical implementation.}
          Incorporating the relaxation kernel into a Boltzmann code would allow
          for global parameter scans and quantitative comparison with
          cosmological data, enabling a full observational assessment of RDCC.

    \item \textbf{Emergent geometry.}
          The relation \(\Gamma(t) \propto H(t)\) suggests that aspects of
          large-scale geometry may be understood in terms of correlation decay
          rather than fundamental spacetime primitives, providing a potential
          bridge between RDCC and information-theoretic approaches to gravity.
\end{itemize}

The results of this Companion provide the microscopic foundation for the
relaxation dynamics used throughout the RDCC framework and set the stage for
further developments in the information-theoretic and phenomenological
structure of cyclic cosmology.

\appendix
\section*{Appendix A: Projective Compactification of the CPT Fixed Point}
\addcontentsline{toc}{section}{Appendix A: Projective Compactification of the CPT Fixed Point}

The global RDCC state lives on the bipartite Hilbert space \(\Hglobal\), with the
two sectors related by global CPT symmetry. At the CPT-invariant bounce, the
coherence parameter satisfies \(\alpha(t_b) \to 0\), the two sectors become
maximally entangled, and the coarse-grained entropy of the global state reaches
its minimum. This hypersurface is the structural fixed point of the CPT
symmetry. In this appendix we formalize this fixed point using a projective
compactification of the time coordinate.

\subsection*{Projective Structure}

We introduce the real projective line \(\mathbb{RP}^{1}\), the one-point
compactification of \(\mathbb{R}\) in which \(+\infty\) and \(-\infty\) are
identified. Points on \(\mathbb{RP}^{1}\) are equivalence classes \([x:y]\) with
\([x:y] = [\lambda x : \lambda y]\) for \(\lambda \neq 0\). The ordinary reals
embed via \(t \mapsto [t:1]\), while the projective point at infinity is
\([1:0]\).

To describe the neighbourhood of the bounce, we define a projective deviation
coordinate


\[
\tau = \frac{1}{a - a_b},
\]


with \(a_b = 1\) the minimal scale factor. As \(a \to 1^{+}\) (the \(+\) sector),
one has \(\tau \to +\infty\); as \(a \to 1^{-}\) (the \(-\) sector), one has
\(\tau \to -\infty\). Both limits map to the same projective point:


\[
\lim_{\tau \to \pm\infty} [1 : 1/\tau] = [1 : 0].
\]


Thus the bounce corresponds to the projective point where the two time
directions are identified.

\subsection*{CPT Action}

In these coordinates, the global CPT transformation acts as


\[
\tau \mapsto -\tau,
\]


enforcing


\[
\rplus(+\tau) = \Theta\, \rminus(-\tau)\, \Theta^{-1},
\]


where \(\Theta\) denotes the CPT operator. At the fixed point \(\tau = 0\), the
two reduced density matrices coincide and the entanglement entropy is maximal,
consistent with \(\alpha(t_b) = 0\).

\subsection*{Geometric Interpretation of Relaxation}

The multiplicative evolution of the coherence parameter,


\[
\RelaxEq,
\qquad
\Gamma(t) \propto H(t),
\]


describes the flow of \(\alpha(a)\) away from the projective fixed point. As the
Universe moves along the compactified line, the distance from \(\tau = 0\)
increases and the two sectors decohere at a rate set by the Hubble expansion.
The projective compactification therefore provides a geometric interpretation of
the relaxation dynamics: the bounce is a removable singularity in the global
state, and the two thermodynamic arrows emerge as opposite directions along the
projective line.

\subsection*{Coherence Length}

The characteristic coherence length


\[
\lcoh \sim \frac{c}{H(t)}
\]


emerges naturally as the scale controlling how rapidly correlations can propagate
across the projective fixed point before relaxation enforces sectoral
projection. This geometric picture complements the dynamical and microscopic
derivations presented in the main text and clarifies the structural unity of the
RDCC framework.


\end{document}
